\section{Contribución Audiovisual}
\label{sec:contribucion_av}

La contribución audiovisual designa el transporte de la señal de vídeo y audio desde
las fuentes de captura, cámaras, ordenadores de presentación y reproductores, hasta
el mezclador o núcleo del sistema de producción. Es el primer eslabón de la cadena y
su elección determina la calidad máxima disponible, la latencia de extremo a extremo
y el coste de la infraestructura. En esta sección se describen las principales
interfaces y protocolos empleados: SDI, SMPTE ST~2022-6, SMPTE ST~2110, NDI y SRT.

\subsection{Interfaces y protocolos}

\begin{itemize}
  \item \textbf{SDI (\textit{Serial Digital Interface}):} estándar histórico de la
  industria broadcast para el transporte de vídeo sin comprimir sobre cable coaxial o
  fibra óptica~\cite{smpte_sdi}. Ofrece latencia sub-imagen y sincronismo garantizado,
  pero está limitado a distancias de hasta 100~m sin regeneradores y requiere
  infraestructura especializada de alto coste.

  \item \textbf{SMPTE ST~2022-6 (SDI sobre IP):} encapsula la señal SDI íntegra en
  paquetes UDP/IP, facilitando la migración de infraestructuras SDI existentes hacia
  redes Ethernet estándar~\cite{smpte_st2022_6}. Incorpora corrección de errores FEC
  para compensar las pérdidas propias de las redes IP.

  \item \textbf{SMPTE ST~2110 (Producción IP nativa):} estándar publicado en 2017 que
  transporta vídeo, audio y metadatos como \textbf{flujos IP independientes},
  proporcionando mayor flexibilidad de encaminamiento que SDI~\cite{smpte_st2110_20}.
  La sincronización entre flujos se consigue mediante \textbf{PTP} (IEEE~1588)~\cite{ieee1588_ptp},
  con precisión de microsegundos. Logra ahorros de ancho de banda del 15--30\% respecto
  a ST~2022-6 y está muy extendido en nuevas instalaciones broadcast profesionales.

  \item \textbf{NDI (\textit{Network Device Interface}):} protocolo propietario de
  Vizrt/NewTek~\cite{ndi_vizrt} que transporta vídeo con compresión ligera sobre redes
  LAN Ethernet con latencia de un fotograma. Ampliamente soportado por cámaras PTZ,
  mezcladores hardware y herramientas de software de producción.

  \item \textbf{SRT (\textit{Secure Reliable Transport}) en contribución:} permite
  transportar la señal desde el lugar del evento hasta el centro de producción remoto
  sobre Internet pública con latencia inferior a 1~s y corrección de errores
  adaptativa~\cite{sharabayko2024_srt}. FFmpeg lo soporta como entrada y salida, lo
  que permite integrarlo directamente en pipelines Voctomix.
\end{itemize}

\begin{table}[H]
  \centering
  \caption{Comparativa de interfaces y protocolos de contribución audiovisual.}
  \label{tab:contrib_interfaces}
  \begin{tabular}{|l|l|c|c|l|}
    \hline
    \textbf{Estándar} & \textbf{Medio} & \textbf{Velocidad máx.} & \textbf{Latencia} & \textbf{Uso típico} \\
    \hline
    SDI (12G)       & Coaxial/fibra & 11{,}88 Gbit/s & Sub-imagen  & Estudio broadcast \\
    SMPTE ST 2022-6 & Ethernet IP   & 3 Gbit/s       & Sub-imagen  & Migración SDI/IP \\
    SMPTE ST 2110   & Ethernet IP   & Variable       & Sub-imagen  & Producción IP nativa \\
    NDI             & Ethernet LAN  & $\sim$100 Mbit/s & 1 fotograma & Software, PTZ \\
    SRT             & Internet/LAN  & Variable       & $<$1~s      & Contribución REMI IP \\
    \hline
  \end{tabular}
\end{table}

\subsection{Códecs de contribución}

En la etapa de contribución se prioriza la calidad y la ausencia de pérdidas sobre la
eficiencia de compresión, a diferencia de la distribución al espectador final:

\begin{itemize}
  \item \textbf{Vídeo sin comprimir (RAW):} el vídeo se transporta sin ningún tipo de
  compresión, preservando toda la información de la señal original. Los formatos más
  utilizados son UYVY (crominancia 4:2:2 empaquetada, 2~bytes/píxel) e I420 (YUV
  4:2:0 planar). Voctomix emplea I420 internamente para todo el procesamiento mediante
  el motor GStreamer, transportado sobre conexiones TCP en contenedor Matroska.

  \item \textbf{Apple ProRes}~\cite{apple_prores}: familia de códecs de alta calidad
  desarrollados por Apple, diseñados para producción profesional. Utilizan
  exclusivamente compresión intra-frame (cada fotograma es independiente), lo que
  minimiza la degradación acumulativa en múltiples etapas de procesamiento. ProRes~422
  HQ es el estándar en producción televisiva para postproducción.

  \item \textbf{Avid DNxHD/DNxHR}~\cite{avid_dnxhd}: familia equivalente a ProRes
  desarrollada por Avid, también basada en compresión intra-frame. Ampliamente
  utilizada en flujos de trabajo de postproducción broadcast y en integración con
  sistemas de edición no lineal Avid Media Composer.

  \item \textbf{JPEG~2000}~\cite{jpeg2000_iso}: códec que ofrece compresión sin
  pérdidas (\textit{lossless}) o con pérdidas mínimas y latencia muy reducida gracias
  a su arquitectura de codificación por \textit{tiles} (bloques) independientes. Utilizado en
  sistemas de contribución broadcast de alta calidad y en señalización digital.
\end{itemize}

\subsection{Formatos y protocolos de transporte}
\label{sec:contrib_formatos_protocolos}

El formato contenedor define cómo se encapsulan y multiplexan los flujos de vídeo y
audio para su transporte entre dispositivos:

\subsubsection*{Formatos}

\begin{itemize}
  \item \textbf{Matroska (MKV)}~\cite{matroska_spec}: contenedor flexible y extensible
  basado en un formato binario abierto. Voctomix lo emplea como protocolo de transporte
  interno: las fuentes envían señal I420 y audio PCM~S16LE a 48~kHz encapsulados en MKV
  sobre conexiones TCP, lo que garantiza la sincronía audio-vídeo y permite multiplexar
  múltiples pistas en un único flujo sin overhead significativo.

  \item \textbf{MPEG-TS (\textit{Transport Stream}, ISO~13818-1)}~\cite{iso_mpegts}:
  estándar de broadcast diseñado para transmisión con posibilidad de pérdidas. Su
  robustez frente a errores lo hace idóneo para contribuciones sobre redes no fiables
  y para la integración con sistemas de distribución broadcast.

  \item \textbf{MP4 / ISOBMFF (\textit{ISO Base Media File Format})}~\cite{iso_isobmff}:
  contenedor de referencia para distribución por HTTP progresivo y streaming adaptativo.
  También utilizado como formato de grabación en cámaras profesionales y sistemas de
  contribución modernos.
\end{itemize}

\subsubsection*{Protocolos}

\begin{itemize}
  \item \textbf{RTP (\textit{Real-time Transport Protocol})}~\cite{rfc3550_rtp}:
  protocolo de transporte en tiempo real definido en la RFC~3550, utilizado junto con
  RTSP para la distribución de flujos multimedia. Habitual en cámaras IP profesionales,
  sistemas SMPTE~2110 y en la transmisión de vídeo sin comprimir sobre redes Ethernet.

  \item \textbf{RTSP (\textit{Real Time Streaming Protocol})}~\cite{rfc7826_rtsp}:
  protocolo de señalización definido en la RFC~7826 que se emplea para establecer y
  controlar sesiones de streaming multimedia. Actúa como ``mando a distancia'' del
  flujo RTP: permite iniciar, pausar y detener la transmisión de una fuente IP sin
  interrumpir la conexión.

  \item \textbf{SRT (\textit{Secure Reliable Transport})}~\cite{sharabayko2024_srt}:
  protocolo de código abierto desarrollado por Haivision con corrección de errores ARQ
  y cifrado AES-128/256. En el contexto de contribución se utiliza para transportar la
  señal desde el lugar del evento al centro de producción remoto sobre redes no fiables
  (véase Sección~\ref{sec:distribucion_av} para su rol en distribución).
\end{itemize}
